\section{Interpretované jazyky}\label{sec:interpretovane-jazyky}

Dosud jsme předpokládali,~že překladač vytvoří cílový program a~ten se spustí až potom. Druhou základní možností je vykonávat program prostřednictvím jiného programu --- interpretu.

\emph{Interpret} načítá reprezentaci programu a~vykonává její příkazy nebo instrukce. Nemusí předem vytvořit samostatný spustitelný soubor se strojovým kódem.

\begin{figure}[h]
    \centering
    \begin{tikzpicture}[
    box/.style={task,minimum width=25mm,minimum height=10mm,font=\small}
]
    \node[font=\bfseries] at (4.95,2.2) {kompilovaný překlad};
    \node[box,fill=blue!10] (src1) at (0,1.1) {zdrojový\\program};
    \node[box,fill=orange!12] (comp) at (3.3,1.1) {překladač};
    \node[box,fill=green!10] (target) at (6.6,1.1) {cílový\\program};
    \node[box,fill=green!10] (result1) at (9.9,1.1) {výsledky};
    \node[box,fill=gray!10] (data1) at (6.6,-0.2) {data};
    \draw[diredge] (src1)--(comp);
    \draw[diredge] (comp)--(target);
    \draw[diredge] (data1)--(target);
    \draw[diredge] (target)--(result1);

    \node[font=\bfseries] at (4.95,-1.4) {interpretovaný překlad};
    \node[box,fill=blue!10] (src2) at (0,-2.5) {zdrojový\\program};
    \node[box,fill=orange!12] (interp) at (4.95,-2.5) {interpret};
    \node[box,fill=green!10] (result2) at (9.9,-2.5) {výsledky};
    \node[box,fill=gray!10] (data2) at (4.95,-3.8) {data};
    \draw[diredge] (src2)--(interp);
    \draw[diredge] (data2)--(interp);
    \draw[diredge] (interp)--(result2);
\end{tikzpicture}

    \caption{Zjednodušené porovnání kompilovaného a~interpretovaného překladu.}
    \label{fig:kompilace-interpretace}
\end{figure}

\warningbox{Přesnější je mluvit o~kompilované nebo interpretované \emph{implementaci},~nikoliv o~neměnné vlastnosti jazyka. Tentýž jazyk může mít klasický překladač,~interpret i~běhové prostředí kombinující více přístupů.}

\paragraph{Výhody a~nevýhody interpretace.}

Interpretované vykonávání usnadňuje přenos programu mezi platformami: stačí,~aby na nich existoval odpovídající interpret. Často také zjednodušuje interaktivní práci,~ladění a~rychlé úpravy programu.

Za běhu se však opakuje práce,~kterou by klasický překladač mohl provést předem. Interpret musí načítat instrukce své vnitřní reprezentace,~rozhodovat o~jejich významu a~spravovat zásobník či prostředí proměnných. Čistá interpretace proto bývá pomalejší než přímé vykonávání dobře přeloženého strojového kódu. Navíc musí být při spuštění přítomno běhové prostředí.

\subsection{Virtuální stroj}

Náš překladač vytváří bajtkód. Program \ref{lst:vm} jej vykonává pomocí zásobníku a~slovníku proměnných. Jde tedy o~malý \emph{virtuální stroj}: instrukce nejsou instrukcemi skutečného procesoru,~ale mají přesně definovaný účinek v~našem programu.

Virtuální stroj odděluje zdrojový jazyk od skutečného procesoru. Překladač stačí přizpůsobit jedné sadě virtuálních instrukcí a~pro každou podporovanou platformu vytvořit jejich interpret. Tentýž bajtkód pak lze spustit na různém hardwaru,~pokud na něm existuje odpovídající běhové prostředí.

\lstinputlisting[
    caption={Interpret zásobníkového bajtkódu.},
    label={lst:vm}
]{04-kompilatory/code/mini_compiler/virtual_machine.py}

U~binární operace leží na vrcholu zásobníku pravý operand a~pod ním levý. Virtuální stroj je odebere,~spočítá výsledek a~ten opět vloží na zásobník. Instrukce \texttt{HALT} ukončí interpretační smyčku.

Zásobníkové instrukce nemusejí přímo uvádět registry operandů,~a~proto mají jednoduchý formát. Cenou je větší počet operací se zásobníkem. Jiné virtuální stroje mohou být registrové; volba reprezentace ovlivňuje velikost bajtkódu,~rychlost interpretace i~složitost generátoru.

\subsection{JIT kompilace}

\emph{JIT kompilace} (z~anglického \emph{just-in-time compilation}) překládá mezikód do strojového kódu až za běhu programu. Běhové prostředí může nejprve bajtkód interpretovat,~sledovat často používané části a~jen ty později přeložit a~optimalizovat.

\begin{figure}[h]
    \centering
    \resizebox{\textwidth}{!}{\begin{tikzpicture}[
    box/.style={task,minimum width=29mm,font=\small}
]
    \node[box,fill=blue!10] (source) at (0,0) {zdrojový\\kód};
    \node[box,fill=orange!12] (compiler) at (3.3,0) {kompilátor};
    \node[box,fill=green!10] (bytecode) at (6.6,0) {bajtkód /\\mezikód};
    \node[box,fill=orange!12] (runtime) at (9.9,0) {běhové\\prostředí};
    \node[box,fill=green!10] (machine) at (13.2,0) {strojový\\kód};

    \draw[diredge] (source)--(compiler);
    \draw[diredge] (compiler)--(bytecode);
    \draw[diredge] (bytecode)--(runtime);
    \draw[diredge] (runtime)--(machine);
    \node[below=4mm of runtime,align=center,font=\small]
        {interpretace a případně\\JIT kompilace často používaných částí};
\end{tikzpicture}
}
    \caption{Překlad přes mezikód s~možnou JIT kompilací.}
    \label{fig:jit}
\end{figure}

Tento přístup kombinuje přenositelnost mezikódu s~rychlejším vykonáváním často používaných částí. JIT navíc zná skutečné chování programu za běhu,~a~může tedy využít informace,~které statický překladač neměl. Platí za to časem a~pamětí spotřebovanými na překlad během spuštění. Není proto automaticky nejlepší pro každý program. Známými příklady běhových prostředí využívajících JIT jsou virtuální stroj Javy,~prostředí .NET a~moderní implementace JavaScriptu.

\subsubsection*{Celý překlad}

Program \ref{lst:main_compiler} spojuje všechny vytvořené části. Zdrojový text projde lexerem,~parserem,~sémantickou kontrolou,~výpočtem konstant a~generátorem bajtkódu. Nakonec jej vykoná virtuální stroj.

Každá fáze přijímá výsledek předchozí fáze a~vrací jednoznačně určenou reprezentaci. Pokud nastane chyba,~má být oznámena tam,~kde ji lze nejlépe vysvětlit: neplatný znak v~lexeru,~chybějící část příkazu v~parseru a~neznámé jméno při sémantické kontrole. Tím se chyba nepropaguje do generátoru kódu jako obtížně srozumitelný problém.

\lstinputlisting[
    caption={Sestavení a~spuštění celého malého překladače.},
    label={lst:main_compiler}
]{04-kompilatory/code/mini_compiler/main.py}

Pro přiložený vstup optimalizátor nahradí výraz $2+3\cdot4$ číslem $14$. Vygenerovaný bajtkód nejprve uloží hodnotu do proměnné \texttt{x},~vypíše ji,~spočítá proměnnou \texttt{y} a~vypíše také její hodnotu. Výstupem je
\begin{verbatim}
14.0
4.0
\end{verbatim}

Ukázku lze použít i~jako test celého překladového řetězce. Změníme-li vstup,~můžeme porovnat očekávaný výstup s~výsledkem virtuálního stroje; samostatnými testy navíc ověříme tokeny,~strom a~bajtkód. Takové rozdělení testů pomáhá určit,~ve které fázi případná chyba vznikla.

\tipbox{Ukázkový překladač je malý,~ale hranice mezi částmi odpovídají skutečným systémům: lexer neřeší gramatiku,~parser neřeší deklarace,~sémantická analýza nemění význam programu a~generátor kódu už dostává ověřenou reprezentaci. Právě toto rozdělení umožňuje každou část samostatně testovat a~později nahradit propracovanější variantou.}
